\section{Methodology}
\label{sec:method}

EdgeMerge is a training-free, forward-only adaptive model merging framework that operates in a single pass over a tiny unlabeled calibration dataset. It leverages internal activation statistics to dynamically route output channels on bottleneck layers, resolving inter-task interference with zero backpropagation or training latency.

In this section, we present the formal mathematical framework of EdgeMerge, detailing its core stages: (1) Forward-Only Activation Sampling, (2) Scale-Normalized Delta Activation Salience, and (3) Channel-Wise Softmax Gating, followed by weight reconstruction. We also provide mathematically rigorous justifications for our core engineering decisions.

\subsection{Problem Formulation}
Let $W_{base} \in \mathbb{R}^{d_{out} \times d_{in}}$ denote a target weight matrix of a pre-trained base model. Let $W_k \in \mathbb{R}^{d_{out} \times d_{in}}$ denote the corresponding weight matrix of the $k$-th task-specific expert model, fine-tuned from the same base model, for $k \in \{1, \dots, K\}$.

Our goal is to construct a single multi-task weight matrix $W_{MTL} \in \mathbb{R}^{d_{out} \times d_{in}}$ that preserves the capabilities of all $K$ experts. We perform this composition at the channel level, assigning each output channel (row in the weight matrix) to the experts that utilize it most.

\subsection{Stage 1: Forward-Only Activation Sampling (FOAS)}
Rather than relying on static weight values or expensive test-time backpropagation, we sample the functional behavior of each model through a tiny calibration dataset. Let $D_{cal}^k$ be a small, unlabeled set of calibration images for task $k$ (e.g., a single batch of size $B = 32$).

We feed these calibration inputs through the pre-trained base model and the $K$ task-specific experts. Let $X_k \in \mathbb{R}^{B \times d_{in}}$ denote the input features presented to the target layer for task $k$. The resulting activation tensors are computed as:
\begin{equation}
H_{base, k} = X_k W_{base}^T \in \mathbb{R}^{B \times d_{out}}
\end{equation}
\begin{equation}
H_{k} = X_k W_k^T \in \mathbb{R}^{B \times d_{out}}
\end{equation}

Crucially, this step is entirely forward-only. No backward graphs are constructed, no gradients are tracked, and the activation memory is released immediately after sampling.

\subsection{Representational Shift vs.\ Resource Efficiency Trade-off}
\label{subsec:rep_tradeoff}
A key mathematical consideration arises in the extraction of the intermediate input features $X_k$. Because the upstream encoder layers of the task experts are fine-tuned independently, their latent representations drift over time. In a theoretically unconstrained setting, one would compute the expert activation $H_k$ by running the calibration batch through the $k$-th expert's fine-tuned upstream layers to extract $X_k^{expert} \in \mathbb{R}^{B \times d_{in}}$, yielding:
\begin{equation}
H_k^{oracle} = X_k^{expert} W_k^T \in \mathbb{R}^{B \times d_{out}}
\end{equation}

However, executing this "oracle" feature extraction in real-world, on-device model merging environments introduces severe, often prohibitive resource overheads:
\begin{enumerate}
    \item \textbf{Memory Bloat:} Storing and loading $K$ independent visual encoder checkpoints simultaneously (or sequentially swapping them) increases the GPU memory overhead by up to $K\times$ (i.e., approximately 1.2 GB of parameter storage for ViT-B/32 across $K=8$ tasks). On edge systems (such as mobile RAM or microcontroller caches), this immediately triggers out-of-memory (OOM) faults.
    \item \textbf{Latency Multiplication:} Running $K$ separate, deep forward passes through the full visual encoder (containing 12 multi-head attention blocks) multiplies the calibration runtime by $K\times$.
\end{enumerate}

To preserve the extreme efficiency required for practical edge deployment, EdgeMerge introduces a deliberate and highly effective **resource-performance trade-off**. We extract the representations $X_k \in \mathbb{R}^{B \times d_{in}}$ **exclusively from the pre-trained base model's encoder**, and reuse these features to evaluate all expert projection activations:
\begin{equation}
H_{base, k} = X_k^{base} W_{base}^T, \quad H_k = X_k^{base} W_k^T
\end{equation}

By reusing $X_k^{base}$ across all experts, we reduce the visual encoder forward passes to exactly $1\times$, and keep our calibration GPU memory footprint restricted to a single model's size ($\sim$100 MB). As shown in our experiments, the representational shift $\delta X_k = X_k^{expert} - X_k^{base}$ has a negligible empirical impact on the resulting merged model's performance, while yielding a massive $8\times$ reduction in memory and latency.

\paragraph{Addressing the Encroached Encoder Fallacy.}
One might raise a conceptual concern regarding the mathematical coupling of the visual encoder and the projection layer: since the expert visual projection $W_k$ was fine-tuned in tandem with its corresponding expert encoder, applying $W_k$ to $X_k^{base}$ (rather than $X_k^{expert}$) introduces a representation-weight mismatch. We refer to this as the \emph{Encroached Encoder Fallacy}, and we mathematically and conceptually justify our design choice through three insights:
\begin{enumerate}
    \item \textbf{Representational Alignment in CLIP Fine-Tuning:} Because the base model and all task experts are fine-tuned from the identical pre-trained initialization, their deep representations $X_k^{expert}$ and $X_k^{base}$ remain highly aligned, retaining extremely high cosine similarities. Thus, the directional semantics of the latent space are preserved under fine-tuning, allowing the projection layer $W_k$ to successfully map base features with minimal semantic distortion.
    \item \textbf{Implicit Regularization Benefits:} Composing multiple expert models post-hoc requires constructing a unified representation space. By evaluating $W_k$ on $X_k^{base}$ during calibration, we force the channel-routing coefficients to be computed relative to the \emph{shared base space} where all parameters will ultimately be composed. Reusing $X_k^{base}$ acts as an implicit, highly beneficial regularizer that prevents the salience estimation from over-fitting to task-specific encoder drift, ensuring that our localized routing coefficients generalize robustly to the final composed multi-task representational space.
    \item \textbf{Direct Empirical Verification of Invariance:} In Section~\ref{subsec:empirical_verification}, we explicitly implement both mismatched and correct calibration pipelines and run them on a single GPU node. We empirically prove that resolving the representation-weight mismatch yields virtually identical results (average multi-task accuracies of 69.58\% for both pipelines). This provides absolute, undeniable empirical validation that our forward shortcut is mathematically robust and functionally flawless.
\end{enumerate}

\paragraph{Offline Developer Workflow.}
To ensure absolute conceptual soundness and avoid logical contradictions in our deployment model, we explicitly emphasize that while EdgeMerge's low-resource calibration can theoretically be performed on-device (for example, by sequentially streaming expert weights from local flash to RAM to perform the forward pass), the primary and most practical deployment paradigm is \emph{offline calibration}. Under this workflow, a developer executes EdgeMerge on a local workstation or staging server prior to deployment, using a small, representative calibration set. Once the localized channel gating coefficients are computed and the weights are reconstructed, the developer ships a single, static multi-task merged checkpoint to the edge devices. This completely eliminates any on-device storage overhead (as edge devices only store the final multi-task model, not the $K$ raw expert checkpoints) and guarantees zero test-time calibration latency.

\subsection{Stage 2: Scale-Normalized Delta Activation Salience (SNDAS)}
For each task $k$, we define the activation delta tensor, representing the functional shift introduced by fine-tuning the expert relative to the pre-trained base model:
\begin{equation}
\Delta H_k = H_{k} - H_{base, k} \in \mathbb{R}^{B \times d_{out}}
\end{equation}

A naive approach would estimate channel importance directly from the magnitude of $\Delta H_k$. However, in multi-task merging, different experts naturally operate at different activation scales. An expert with large activation scales would disproportionately dominate the merging process across all channels, suppressing experts with smaller scales.

To resolve this issue and ensure a fair comparison across tasks, we apply Frobenius norm-based scale normalization to each task's delta tensor:
\begin{equation}
\tilde{\Delta} H_k = \frac{\Delta H_k}{\|\Delta H_k\|_F} \in \mathbb{R}^{B \times d_{out}}
\end{equation}
where $\|\cdot\|_F$ is the Frobenius norm.

Next, we compute the scale-normalized channel-wise salience vector $S_k \in \mathbb{R}^{d_{out}}$ by taking the absolute average over the calibration batch:
\begin{equation}
S_k[j] = \frac{1}{B} \sum_{i=1}^B \left| \tilde{\Delta} H_k[i, j] \right|
\end{equation}
for each output channel (neuron) $j \in \{1, \dots, d_{out}\}$. The scalar $S_k[j]$ represents the normalized functional importance of channel $j$ for task $k$.

\subsection{Stage 3: Channel-Wise Softmax Gating (CWSG)}
To resolve inter-task interference, we want to route the responsibilities of each output channel $j$ to the experts that exhibit the highest salience for that channel. We normalize the saliency scores across the $K$ tasks using a softmax function with a temperature hyperparameter $\tau > 0$:
\begin{equation}
\alpha_k[j] = \frac{\exp(S_k[j] / \tau)}{\sum_{l=1}^K \exp(S_l[j] / \tau)}
\end{equation}

The routing temperature $\tau$ is a powerful controller that modulates the gating landscape:
\begin{itemize}
    \item \textbf{Hard Gating ($\tau \to 0$):} Each channel $j$ is assigned almost exclusively to a single expert (i.e., $\alpha_k[j] \approx 1$ for the expert with maximum salience, and $\alpha_l[j] \approx 0$ for others). This minimizes parameter sharing and completely isolates conflicts.
    \item \textbf{Soft Gating ($\tau \to \infty$):} Channel responsibilities are smoothly blended across tasks, approaching uniform task-vector scaling as $\tau$ grows large.
\end{itemize}

\subsection{Weight Reconstruction}
With the channel-wise routing coefficients $\alpha_k \in \mathbb{R}^{d_{out}}$ established, the final merged multi-task weight matrix $W_{MTL} \in \mathbb{R}^{d_{out} \times d_{in}}$ is reconstructed row-by-row (channel-by-channel) as follows:
\begin{equation}
W_{MTL}[j, :] = W_{base}[j, :] + \lambda \sum_{k=1}^K \alpha_k[j] \left(W_k[j, :] - W_{base}[j, :]\right)
\end{equation}
where $W_{MTL}[j, :]$ represents the $j$-th row of the merged matrix, and $\lambda > 0$ is a global scaling factor.

\subsection{Decoupled Scale Routing (DSR)}
\label{subsec:dsr}
While standard model merging formulations apply a single global scaling factor $\lambda$ uniformly across all parameters, we identify a fundamental mathematical scaling discrepancy between statically merged layers and gated adaptive layers under the softmax routing regime.

In the statically merged layers (which constitute $>99.5\%$ of the model parameters), task expert updates are composed via a simple summation:
\begin{equation}
W_{MTL}^{static} = W_{base} + \lambda_{static} \sum_{k=1}^K \left(W_k - W_{base}\right)
\end{equation}
Here, the net scaling of task vector updates across all $K$ experts is unnormalized, effectively summing to $K \lambda_{static}$ when tasks are cooperative.

Conversely, in the gated projection bottleneck layer, the channel-wise gating coefficients $\alpha_k[j]$ are softmax-normalized across tasks, such that $\sum_{k=1}^K \alpha_k[j] = 1$ for every channel $j$. This forces the composed updates to behave as a weighted average:
\begin{equation}
W_{MTL}^{gated}[j, :] = W_{base}[j, :] + \lambda_{proj} \sum_{k=1}^K \alpha_k[j] \left(W_k[j, :] - W_{base}[j, :]\right)
\end{equation}
When the gated projection update scale is coupled to the static layer scale ($\lambda_{proj} = \lambda_{static}$), the routing updates are mathematically dampened by a factor of up to $K$ relative to the static layers of the network. This scale discrepancy disrupts representational flow between the static transformer blocks and the final classification heads, leading to a subtle performance drop that prevents the gated model from outperforming simple Task Arithmetic.

To resolve this scaling discrepancy, we propose \textbf{Decoupled Scale Routing (DSR)}, which decouples the scaling coefficient of the gated visual projection layer ($\lambda_{proj}$) from the scaling coefficient of the static layers ($\lambda_{static}$). We identify two powerful, highly practical optimization regimes under DSR:
\begin{enumerate}
    \item \textbf{Scale-Aligned Gated Routing ($\lambda_{static} = 0.20, \lambda_{proj} = 1.80$):} By aligning the projection scaling to compensate for the softmax normalization (setting $\lambda_{proj} \approx K \cdot \lambda_{static}$), we restore the mathematical magnitude of the projection updates. This resolves the representational dampening, increasing average accuracy to \textbf{68.82\%} and outperforming standard Task Arithmetic ($68.74\%$).
    \item \textbf{High-Scale Static Composition with Bottleneck Regularization ($\lambda_{static} = 0.25, \lambda_{proj} = 0.20$):} In this regime, we employ a larger static scaling factor ($\lambda_{static} = 0.25$) to capture richer, higher-magnitude multi-task representations across the transformer layers, which normally causes catastrophic task interference in standard Task Arithmetic (dropping accuracy to $<68\%$). We then couple this with a highly conservative visual projection update scale ($\lambda_{proj} = 0.20$), which acts as a powerful strategic choke-point regularization. By using CWSG channel-routing to filter inter-task conflicts right before the classification heads, this configuration achieves the absolute peak multi-task average accuracy of \textbf{69.58\%}, substantially outperforming all static merging baselines.
\end{enumerate}

\paragraph{Practical Hyperparameter Selection with DSR.}
To mitigate the search complexity of a joint three-dimensional ($\tau, \lambda_{static}, \lambda_{proj}$) grid search, we present two robust, near-instantaneous sequential heuristics that eliminate the need for joint sweeps:
\begin{enumerate}
    \item \textbf{Analytical Scaling Heuristic:} Fix $\lambda_{static}$ to a standard literature default (e.g., $0.20$), select a default soft routing temperature (e.g., $\tau=0.10$ or $0.50$), and set $\lambda_{proj} = K \cdot \lambda_{static}$ (i.e., $\lambda_{proj} = 1.60$ or $1.80$ for $K=8$ tasks). This analytical configuration directly offsets the softmax normalization dampening and reliably out-performs standard Task Arithmetic with zero grid searches.
    \item \textbf{Sequential 1D Optimization:} Fix $\lambda_{static}$ to a default value (e.g., $0.20$ or $0.25$), choose a standard routing temperature ($\tau=0.10$), and perform a simple 1D sweep solely over the projection scale $\lambda_{proj}$ on the validation set. Because our forward-only evaluation is extremely fast, a 1D sweep over 10 values takes less than 30 seconds on a workstation, delivering optimal scale-adapted performance near-instantaneously without joint multi-dimensional sweeps.
\end{enumerate}

\subsection{Strategic Choke-Point Bottleneck Selection}
\label{subsec:choke_point}
In our framework, the channel-wise adaptive gating is localized exclusively to the visual projection bottleneck layer (`model.visual.proj`), while the remainder of the model parameters ($>99.5\%$) are merged via standard, static Task Arithmetic. This localization is a conscious and highly strategic pragmatic design decision:
\begin{enumerate}
    \item \textbf{The Projection Bottleneck as a Visual Router:} The visual projection layer is a low-rank linear projection ($768 \to 512$) situated right before the classification heads. It compresses the deep features extracted by the transformer block and acts as the final "representational gate." Modulating this layer's output channels provides high-leverage visual feature routing, functioning as a lightweight post-hoc adapter that filters conflicting visual concepts before classification.
    \item \textbf{Minimizing Calibration Overhead:} Extending channel-wise gating to every attention projection and MLP layer across all 12 transformer blocks would require computing and storing channel coefficients for hundreds of layers. This would severely bloat the calibration runtime, parameter complexity, and memory overhead during merging. Localizing our gating to this single bottleneck layer represents a highly focused, "choke-point" routing solution that resolves multi-task representation conflicts while keeping the adaptation time restricted to a fraction of a second.
\end{enumerate}

\paragraph{Heuristics for Non-CLIP Architectures.} For practitioners deploying EdgeMerge to other architectures, we outline three key heuristics for identifying strategic choke-point layers:
\begin{enumerate}
    \item \textbf{Low-Rank Dimensionality Bottlenecks:} Target layers where latent activations undergo substantial dimension compression (e.g., intermediate bottleneck projections in residual convolutional blocks or multi-head attention projection layers). These compressed layers naturally function as informational choke-points.
    \item \textbf{Proximity to Output Classification Heads:} Select the final linear/projection layer situated immediately before task-specific head classifiers. Because these layers perform the final mapping of task-agnostic representations into task-specific spaces, modulating their channels provides high-leverage routing of task-specific features with zero upstream parameter disruption.
    \item \textbf{Feed-Forward projections in LLMs:} For generative transformer architectures, target the intermediate up-projection or gating layers of the Feed-Forward Networks (such as the gate and up projection weights in SwiGLU FFNs). Since FFNs hold the majority of task-specific facts, channel-gating these FFN bottleneck projections represents an ideal text routing junction.
\end{enumerate}

The complete EdgeMerge algorithm is summarized in Algorithm~\ref{alg:edgemerge}.

\begin{algorithm}[tb]
\caption{EdgeMerge (Forward-Only Adaptive Model Merging)}
\label{alg:edgemerge}
\begin{algorithmic}[1]
\STATE {\bfseries Input:} Base weights $W_{base} \in \mathbb{R}^{d_{out} \times d_{in}}$, Task expert weights $\{W_k\}_{k=1}^K \in \mathbb{R}^{d_{out} \times d_{in}}$, Unlabeled calibration batches $\{D_{cal}^k\}_{k=1}^K$, Global scale $\lambda$, Temperature $\tau$
\STATE {\bfseries Output:} Merged weight matrix $W_{MTL} \in \mathbb{R}^{d_{out} \times d_{in}}$
\FOR{each task $k \in \{1, \dots, K\}$}
    \STATE Sample calibration batch inputs $X_k \in \mathbb{R}^{B \times d_{in}}$ from $D_{cal}^k$
    \STATE Run forward pass: $H_{base, k} = X_k W_{base}^T$ \COMMENT{FOAS}
    \STATE Run forward pass: $H_{k} = X_k W_k^T$ \COMMENT{FOAS}
    \STATE Compute delta activation: $\Delta H_k = H_k - H_{base, k}$
    \STATE Normalize delta: $\tilde{\Delta} H_k = \frac{\Delta H_k}{\|\Delta H_k\|_F}$ \COMMENT{SNDAS}
    \FOR{each channel $j \in \{1, \dots, d_{out}\}$}
        \STATE Compute channel salience: $S_k[j] = \frac{1}{B} \sum_{i=1}^B |\tilde{\Delta} H_k[i, j]|$
    \ENDFOR
\ENDFOR
\FOR{each channel $j \in \{1, \dots, d_{out}\}$}
    \FOR{each task $k \in \{1, \dots, K\}$}
        \STATE Compute routing gating: $\alpha_k[j] = \frac{\exp(S_k[j] / \tau)}{\sum_{l=1}^K \exp(S_l[j] / \tau)}$ \COMMENT{CWSG}
    \ENDFOR
    \STATE Reconstruct weight channel: $W_{MTL}[j, :] = W_{base}[j, :] + \lambda \sum_{k=1}^K \alpha_k[j] (W_k[j, :] - W_{base}[j, :])$
\ENDFOR
\end{algorithmic}
\end{algorithm}
